\documentclass[10pt, aspectratio=169]{beamer}
\usepackage[utf8]{inputenc}
\usepackage[portuguese]{babel}
\usepackage{amsmath, amssymb, amsfonts}
\usepackage{graphicx}
\usepackage{booktabs}
\usepackage{tikz}
\usetikzlibrary{shapes, arrows.meta, positioning, calc}
\usepackage{listings}
\usepackage{xcolor}

% Color Palette - Professional Data Science & Interactive Class Theme
\definecolor{navyblue}{RGB}{30, 41, 59}       % #1E293B Primary dark
\definecolor{skycyan}{RGB}{14, 165, 233}      % #0EA5E9 Secondary blue/cyan
\definecolor{ambergold}{RGB}{245, 158, 11}    % #F59E0B Accent amber
\definecolor{darkgray}{RGB}{51, 65, 85}       % Dark text / subheaders
\definecolor{lightgray}{RGB}{241, 245, 249}   % Card / block background
\definecolor{codebg}{RGB}{248, 249, 250}      % Code background
\definecolor{codegreen}{RGB}{34, 197, 94}     % Code comments/strings
\definecolor{crimsonred}{RGB}{225, 29, 72}     % Highlight red
\definecolor{livepurple}{RGB}{147, 51, 234}    % Purple for Live Coding

% Beamer Theme Settings
\usetheme{Madrid}
\usecolortheme{whale}

\setbeamercolor{structure}{fg=navyblue}
\setbeamercolor{palette primary}{bg=navyblue,fg=white}
\setbeamercolor{palette secondary}{bg=skycyan,fg=white}
\setbeamercolor{palette tertiary}{bg=darkgray,fg=white}
\setbeamercolor{titlelike}{bg=navyblue,fg=white}

\setbeamercolor{block title}{bg=navyblue,fg=white}
\setbeamercolor{block body}{bg=lightgray,fg=black}

\setbeamercolor{block title alert}{bg=crimsonred,fg=white}
\setbeamercolor{block body alert}{bg=lightgray,fg=black}

\setbeamercolor{block title example}{bg=skycyan,fg=white}
\setbeamercolor{block body example}{bg=lightgray,fg=black}

% Custom Block for Live Coding
\newenvironment{liveblock}[1]{%
    \setbeamercolor{block title}{bg=livepurple,fg=white}%
    \setbeamercolor{block body}{bg=lightgray,fg=black}%
    \begin{block}{#1}%
}{%
    \end{block}%
}

\setbeamertemplate{navigation symbols}{}
\setbeamertemplate{footline}[frame number]

% Python Listing Setup
\lstset{
    language=Python,
    backgroundcolor=\color{codebg},
    basicstyle=\ttfamily\scriptsize\color{navyblue},
    keywordstyle=\bfseries\color{skycyan},
    stringstyle=\color{codegreen},
    commentstyle=\itshape\color{gray},
    numbers=none,
    breaklines=true,
    frame=single,
    rulecolor=\color{lightgray}
}

% Title Slide Metadata
\title[Ciência de Dados: ICs e Hipóteses]{Ciência de Dados -- Aula 09}
\subtitle{Intervalos de Confiança Paramétricos e Testes de Hipóteses\\{\small Abordagens Empíricas via Simulação e Paramétricas via Teorema do Limite Central}}
\author[ADS]{Curso de Análise e Desenvolvimento de Sistemas (ADS)}
\institute[ADS]{Disciplina: Ciência de Dados}
\date{\today}

\begin{document}

% -----------------------------------------------------------------------------
% Slide 1: Title
% -----------------------------------------------------------------------------
\begin{frame}
    \titlepage
\end{frame}

% -----------------------------------------------------------------------------
% Slide 2: Agenda
% -----------------------------------------------------------------------------
\begin{frame}{Agenda da Aula 09}
    \tableofcontents
\end{frame}

% =============================================================================
\section{1. Motivação Real e Conceito de Intervalo de Confiança (IC)}
% =============================================================================

\begin{frame}{1.1 Motivação Real: O Estudo da COVID-19 (UFPel 2020)}
    \begin{block}{O Problema Real de Saúde Pública e Decisão}
        Em maio de 2020, cientistas da Universidade Federal de Pelotas (UFPel) realizaram um estudo nacional de testagem sorológica para estimar a proporção de infectados por COVID-19 no Brasil.
        \begin{itemize}
            \item \textbf{Resultado da Amostra:} Prevalência estimada em \textbf{$1{,}9\%$}.
            \item \textbf{Margem de Erro Comunicada:} Intervalo de Confiança de 95\% entre \textbf{$[1{,}7\%, 2{,}1\%]$}.
        \end{itemize}
    \end{block}

    \vspace{0.2cm}

    \begin{alertblock}{Testando a Hipótese de Imunidade de Rebanho}
        \textbf{Pergunta:} "O Brasil atingiu a imunidade de rebanho (necessário pelo menos $50\%$ de soropositivos)?"
        \begin{itemize}
            \item \textbf{Resposta Explicada:} Não! Mesmo considerando a incerteza amostral, o valor de $50\%$ está extremamente distante do intervalo $[1{,}7\%, 2{,}1\%]$.
            \item \textbf{Conclusão:} Rejeita-se a hipótese de imunidade de rebanho com menos de 5\% de chance de erro.
        \end{itemize}
    \end{alertblock}
\end{frame}

\begin{frame}{1.2 O que é um Intervalo de Confiança (IC)?}
    \begin{block}{Definição Formal de Intervalo de Confiança}
        Um 	\textbf{Intervalo de Confiança de 95\%} para um parâmetro populacional (como a média $\mu$) é um intervalo numérico calculado a partir de dados amostrais tal que:
        \begin{itemize}
            \item Se o experimento de amostragem for repetido 100 vezes nas mesmas condições, 	\textbf{95 desses intervalos conterão a verdadeira média da população ($\mu$)}.
        \end{itemize}
    \end{block}

    \vspace{0.2cm}

    \begin{alertblock}{Desmistificando Erros Comuns de Interpretação}
        \begin{itemize}
            \item \textbf{[INCORRETO]:} "Há 95\% de chance de um indivíduo aleatório estar dentro do IC." (IC não fala de indivíduos brutos!)
            \item \textbf{[INCORRETO]:} "95\% dos dados da amostra estão dentro do IC." (IC não mede dispersão de dados!)
            \item \textbf{[CORRETO]:} "Estimamos que a \textbf{média esperada da população} está entre o limite inferior e o limite superior com 95\% de confiança (5\% de chance de erro de amostragem)."
        \end{itemize}
    \end{alertblock}
\end{frame}

% =============================================================================
\section{2. Abordagem 1: IC e Teste de Hipóteses via Simulação Empírica}
% =============================================================================

\begin{frame}{2.1 Estudo de Caso Explicado: O Teste da Moeda Viesada}
    \begin{block}{O Problema do Experimento}
        Você joga uma moeda para cima 30 vezes. A moeda cai em 	\textbf{cara 22 vezes}.
        \textbf{Pergunta de Teste:} "A moeda é justa (não-viesada) ou possui vício a favor de cara?"
    \end{block}

    \vspace{0.2cm}

    \begin{exampleblock}{Resolução via Abordagem Empírica (Sem Fórmulas Decoradas)}
        \begin{enumerate}
            \item \textbf{Hipótese Nula ($H_0$):} A moeda é justa ($P(\text{Cara}) = 0{,}5$). Sob $H_0$, esperamos $30 \times 0{,}5 = 15$ caras.
            \item \textbf{Simulação Computacional:} Repetir 10.000 vezes o experimento de jogar 30 moedas justas e contar o número de caras.
            \item \textbf{Construção do Intervalo Empírico de 95\%:} Extrair os percentis empíricos $2{,}5\%$ e $97{,}5\%$ das 10.000 simulações.
            \item \textbf{Resultado:} O intervalo de 95\% para uma moeda justa em 30 jogadas é 	\textbf{$[10, 20]$ caras}.
            \item \textbf{Decisão Final:} Como 22 caras está 	\textbf{fora} do intervalo $[10, 20]$, a chance de ver 22 caras ao acaso em uma moeda justa é de apenas $0{,}69\%$ ($p < 0{,}05$). 	\textbf{Rejeitamos a hipótese de moeda justa!}
        \end{enumerate}
    \end{exampleblock}
\end{frame}

\begin{frame}[fragile]{[LIVE CODING] Prática 1: Simulação da Moeda Viesada em Python}
    \begin{liveblock}{[PRÁTICA EM SALA] Execução do Teste Empírico via Percentis}
        Executando a simulação de 10.000 lançamentos no Python:
    \end{liveblock}

\begin{lstlisting}
import numpy as np

# Simulando 10.000 experimentos de 30 moedas sob H0 (p=0.5)
np.random.seed(42)
simulacoes = np.random.binomial(n=30, p=0.5, size=10000)

ic_low = np.percentile(simulacoes, 2.5)   # 10
ic_high = np.percentile(simulacoes, 97.5) # 20
p_valor = np.mean(simulacoes >= 22)       # 0.0069 (0.69%)

print(f"Intervalo Empirico de 95% sob H0: [{ic_low:.0f}, {ic_high:.0f}] caras")
print(f"P-Valor (Probabilidade de N >= 22): {p_valor:.4f}")

if p_valor < 0.05:
    print("Rejeitamos H0! A moeda possui vies estatisticamente significante.")
\end{lstlisting}
\end{frame}

% =============================================================================
\section{3. Abordagem 2: IC Paramétrico via Teorema do Limite Central (TCL)}
% =============================================================================

\begin{frame}{3.1 A Equação do IC Paramétrico e Suas Variáveis}
    \begin{block}{Fórmula Geral do Intervalo de Confiança Paramétrico}
        \[ \text{IC} = \bar{x} \pm z^* \cdot \text{SE} = \bar{x} \pm z^* \cdot \frac{s}{\sqrt{n}} \]
        \textbf{Significado Explicado de Cada Símbolo e Variável:}
        \begin{itemize}
            \item $\bar{x}$ (x-barra): 	\textbf{Média amostral} calculada a partir dos dados (estimador pontual da média populacional $\mu$).
            \item $z^*$: 	\textbf{Valor crítico} da Distribuição Normal Padrão associado ao nível de confiança escolhido.
                  \begin{itemize}
                      \item Para 95\% de confiança $\implies z^* = 1{,}96$.
                      \item Para 98\% de confiança $\implies z^* = 2{,}33$.
                      \item Para 99\% de confiança $\implies z^* = 2{,}58$.
                  \end{itemize}
            \item $s$: \textbf{Desvio padrão amostral} (medida de dispersão dos dados brutos).
            \item $n$: \textbf{Tamanho da amostra} (quantidade de observações coletadas).
            \item $\text{SE} = \frac{s}{\sqrt{n}}$: \textbf{Erro Padrão da Média} (\textbf{Standard Error}), que mede a variabilidade da estimativa amostral.
            \item $\text{Margem de Erro} = z^* \cdot \frac{s}{\sqrt{n}}$: O raio de incerteza adicionado e subtraído da média.
        \end{itemize}
    \end{block}
\end{frame}

\begin{frame}{3.2 Estudo de Caso Explicado: Rendimento Acadêmico (RSG UFMG)}
    \begin{block}{O Problema dos Dados Amostrais}
        Entrevistou-se uma amostra aleatória de \textbf{$n = 50$ alunos} do curso de Computação da UFMG sobre seu Rendimento Semestral Global (RSG).
        \begin{itemize}
            \item \textbf{Média Amostral ($\bar{x}$):} $3{,}20$
            \item \textbf{Desvio Padrão Amostral ($s$):} $1{,}74$
        \end{itemize}
        Qual é a estimativa da média real de RSG de toda a população de alunos com 95\% de confiança?
    \end{block}

    \vspace{0.2cm}

    \begin{exampleblock}{Resolução Passo a Passo}
        \begin{enumerate}
            \item \textbf{Cálculo do Erro Padrão (SE):} $\text{SE} = \frac{s}{\sqrt{n}} = \frac{1{,}74}{\sqrt{50}} = \frac{1{,}74}{7{,}071} \approx 0{,}2461$.
            \item \textbf{Cálculo da Margem de Erro (ME):} $\text{ME} = 1{,}96 \times 0{,}2461 \approx 0{,}4823 \approx 0{,}49$.
            \item \textbf{Construção do Intervalo:} $\text{IC} = (3{,}20 - 0{,}49, \; 3{,}20 + 0{,}49) = (2{,}71, \; 3{,}69)$.
            \item \textbf{Resposta Final Explicada:} Estimamos que a média populacional do RSG dos alunos está entre \textbf{$2{,}71$ e $3{,}69$} com 95\% de confiança.
        \end{enumerate}
    \end{exampleblock}
\end{frame}

\begin{frame}[fragile]{[LIVE CODING] Prática 2: IC Paramétrico do RSG em Python}
    \begin{liveblock}{[PRÁTICA EM SALA] Execução do Cálculo de IC com \texttt{scipy.stats}}
        Calculando os limites e o erro padrão exato do RSG:
    \end{liveblock}

\begin{lstlisting}
from scipy import stats
import numpy as np

n = 50
xbar = 3.2
s = 1.74

# 1. Erro Padrao (SE)
se = s / np.sqrt(n) # 0.2461

# 2. Valor critico z* para 95% (1.96)
z_95 = stats.norm.ppf(0.975)

# 3. Margem de Erro e Intervalo
margem = z_95 * se # 0.4823
ic_low, ic_high = xbar - margem, xbar + margem

print(f"SE: {se:.4f} | Margem de Erro: {margem:.4f}")
print(f"Intervalo de Confianca de 95%: ({ic_low:.2f}, {ic_high:.2f})")
\end{lstlisting}
\end{frame}

% =============================================================================
\section{4. Condições de Validade e Dinâmica do Tamanho do IC}
% =============================================================================

\begin{frame}{4.1 Condições Obrigatórias para Usar o IC Paramétrico}
    \begin{block}{Para que o IC via TCL Funcione, Exigem-se 2 Condições:}
        \begin{enumerate}
            \item \textbf{Independência entre as Observações:}
                  \begin{itemize}
                      \item A amostra deve ser não-viesada.
                      \item Se entrevistarmos apenas alunos de um mesmo laboratório, a amostra é viesada. Deve-se amostrar alunos de diferentes institutos/escolas.
                  \end{itemize}
            \item \textbf{Tamanho de Amostra Mínimo ($n \ge 30$):}
                  \begin{itemize}
                      \item O Teorema do Limite Central garante a distribuição normal das médias apenas para $n$ suficientemente grande ($n \ge 30$).
                      \item Para amostras pequenas ($n < 30$), deve-se utilizar a \textbf{Distribuição $t$ de Student} ou \textbf{Bootstrap}.
                  \end{itemize}
        \end{enumerate}
    \end{block}
\end{frame}

\begin{frame}{4.2 Como Alterar a Largura do Intervalo de Confiança?}
    \begin{block}{Análise dos Fatores da Fórmula: $\text{IC} = \bar{x} \pm z^* \cdot \frac{s}{\sqrt{n}}$}
        Como desenvolvedores e analistas de dados, podemos manipular a precisão do IC:
    \end{block}

    \vspace{0.2cm}

    \begin{columns}[T]
        \begin{column}{0.48\textwidth}
            \begin{alertblock}{Variando o Nível de Confiança ($z^*$)}
                \begin{itemize}
                    \item Aumentar a confiança ($95\% \to 98\% \to 99\%$) aumenta o valor de $z^*$ ($1{,}96 \to 2{,}33 \to 2{,}58$).
                    \item \textbf{Efeito:} O intervalo fica \textbf{MAIS LARGO}.
                    \item \textit{Trade-off:} Maior certeza, porém menor precisão.
                \end{itemize}
            \end{alertblock}
        \end{column}

        \begin{column}{0.48\textwidth}
            \begin{exampleblock}{Variando o Tamanho da Amostra ($n$)}
                \begin{itemize}
                    \item Aumentar $n$ reduz o erro padrão $\text{SE} = s/\sqrt{n}$.
                    \item \textbf{Efeito:} O intervalo fica \textbf{MAIS ESTREITO}.
                    \item \textit{Vantagem:} Maior precisão sem perder a confiança de 95\%!
                \end{itemize}
            \end{exampleblock}
        \end{column}
    \end{columns}
\end{frame}

\begin{frame}[fragile]{[LIVE CODING] Prática 3: Tabela de Sensibilidade de Confiança em Python}
    \begin{liveblock}{[PRÁTICA EM SALA] Comparando Confianças de 90\%, 95\%, 98\% e 99\%}
        Visualizando como o intervalo expande ao aumentar a confiança:
    \end{liveblock}

\begin{lstlisting}
confs = [0.90, 0.95, 0.98, 0.99]
for c in confs:
    z = stats.norm.ppf(1 - (1 - c)/2)
    me = z * se
    print(f"Confianca: {c:.0%} | z*: {z:.3f} | Margem: {me:.3f} | IC: ({xbar-me:.2f}, {xbar+me:.2f})")
\end{lstlisting}

    \begin{table}
        \centering
        \tiny
        \begin{tabular}{@{}llll@{}}
            \toprule
            \textbf{Nível de Confiança} & \textbf{Valor Crítico ($z^*$)} & \textbf{Margem de Erro (ME)} & \textbf{Intervalo de Confiança (RSG)} \\ \midrule
            90\% & 1,645 & 0,405 & $(2{,}80, \; 3{,}60)$ \\
            95\% & 1,960 & 0,482 & $(2{,}72, \; 3{,}68)$ \\
            98\% & 2,326 & 0,572 & $(2{,}63, \; 3{,}77)$ \\
            99\% & 2,576 & 0,634 & $(2{,}57, \; 3{,}83)$ \\ \bottomrule
        \end{tabular}
    \end{table}
\end{frame}

% =============================================================================
\section{5. Desafio Prático em Python para ADS}
% =============================================================================

\begin{frame}[fragile]{[LIVE CODING] Desafio Prático: Calculadora Automática de IC para Latência de API}
    \begin{liveblock}{[DESAFIO EM SALA] Criando uma Função Reutilizável em Python}
        Construa uma função que receba uma série de latências de software e retorne o IC de 95\%:
    \end{liveblock}

\begin{lstlisting}
def calcular_ic_api(dados, confianca=0.95):
    n_elem = len(dados)
    xbar_elem = np.mean(dados)
    s_elem = np.std(dados, ddof=1)
    se_elem = s_elem / np.sqrt(n_elem)
    z_critico = stats.norm.ppf(1 - (1 - confianca)/2)
    margem_elem = z_critico * se_elem
    return xbar_elem, (xbar_elem - margem_elem, xbar_elem + margem_elem), margem_elem

np.random.seed(123)
latencias_api = np.random.normal(loc=120, scale=25, size=60)
media_api, (ic_l, ic_h), me_api = calcular_ic_api(latencias_api)

print(f"Media Amostral da API: {media_api:.2f} ms")
print(f"IC 95% do Tempo Medio: ({ic_l:.2f} ms, {ic_h:.2f} ms) | Margem: +- {me_api:.2f} ms")
\end{lstlisting}
\end{frame}

\begin{frame}{Resumo da Aula 09}
    \begin{block}{Síntese do Aprendizado da Aula 09}
        \begin{enumerate}
            \item \textbf{Intervalo de Confiança (IC):} Faixa que contém o verdadeiro valor do parâmetro populacional ($\mu$) em 95\% das amostragens repetidas.
            \item \textbf{Abordagem Empírica vs Paramétrica:} Simulações via percentis empíricos de 95\% vs Fórmula paramétrica do TCL ($\bar{x} \pm 1{,}96 \cdot \text{SE}$).
            \item \textbf{Erro Padrão ($\text{SE} = s/\sqrt{n}$):} Mede a precisão da média amostral.
            \item \textbf{Tomada de Decisão:} Se o valor da hipótese nula $H_0$ cai fora do IC de 95\%, a hipótese é \textbf{rejeitada} com $p < 0{,}05$.
            \item \textbf{Aumentar Precisão:} Aumentar $n$ estreita o intervalo sem perder o nível de confiança.
        \end{enumerate}
    \end{block}
\end{frame}

\end{document}
